\begin{abstract}
\emph{Objective:} Electrocardiogram classifiers trained on complete recordings meet detached and noisy leads in
clinical use. We asked whether quality-gated lead attention or an abstaining evidential output
improves a classifier already trained with simulated lead failure.
\emph{Methods:} On PTB-XL (\nRecordsKept{} records, five diagnostic superclasses) we trained a quality-gated
lead-attention network with an evidential output, three ablations and five baselines under one training budget, and
tested all of them on 31 identical conditions of simulated lead removal and artifacts, with an oracle quality control and paired bootstrap intervals. \emph{Results:} Training with simulated
lead failure raised the macro-averaged area under the receiver operating characteristic curve of the same residual
network from \aucResCleanKTen{} to \aucResDropKTen{} with ten of twelve leads removed, without loss on complete
recordings. With this training, the four attention variants differed by at most \ablRangeKTen{}, and the full system
was not significantly better than the corruption-trained residual network in any of \bootNRows{} paired comparisons
and was worse in \bootNSigNeg{} artifact conditions. Its evidential uncertainty ranked errors no better than its
predicted probabilities and did not increase as leads were removed (\meanUKZero{} to \meanUKTen) while
discrimination fell from \aucLgEdlKZero{} to \aucLgEdlKTen; abstention beyond the clean-data rate came only from a
rule that counts usable leads. \emph{Conclusion:} In this setting the added components gave no measurable benefit
over corruption training. \emph{Significance:} Methods for missing or noisy leads should be compared with baselines
trained on the same failures, and uncertainty estimates should be tested under the failures they are meant to
detect.
\end{abstract}
